% !TEX root = ../../main.tex
\subsection{作品总结}

本文围绕外卖订单通信场景中的身份暴露、订单通信越权、聊天记录篡改和纠纷追责困难等问题，设计并实现了 FoodShield——外卖订单隐私认证与可信审计系统。系统以“日常匿名、异常可溯源”为核心设计思想，将匿名身份、订单认证、实时通信、消息保护、日志完整性验证和管理员审计机制统一到外卖订单通信流程中，形成了较为完整的安全通信闭环。

在身份保护方面，系统通过 PID 匿名身份机制将用户真实身份与通信身份解耦，使骑手端在日常通信过程中无法直接获取用户真实身份或 PID。在订单认证方面，系统基于 HMAC-SM3 构造订单 Token，将订单号、匿名身份和时间戳等信息进行绑定，保证只有合法订单参与方才能进入对应通信房间。在消息保护方面，系统对通信内容进行 SM4 加密存储，并使用 SM3 生成消息摘要，管理员端默认不展示通信明文，从而降低敏感信息泄露风险。

在可信审计方面，系统将消息摘要组织为 Merkle Tree，并通过 Merkle Root 快照实现通信日志完整性验证。当消息内容、消息数量或消息顺序被篡改时，重新计算得到的 Root 将与历史快照不一致，从而能够检测日志篡改行为。在异常处理方面，系统支持用户端和骑手端基于订单号提交溯源申请，管理员经过权限认证后执行条件溯源，并将关键操作写入审计日志，实现了隐私保护与平台治理之间的平衡。

通过系统实现与功能展示可以看出，FoodShield 不只是一个普通的订单聊天系统，而是将密码学机制与外卖订单业务流程相结合的隐私认证与可信审计原型系统。系统实现了用户端、骑手端和管理员端的多端协同，覆盖了注册、下单、认证、通信、审计、验证和溯源等完整流程，具有较好的工程完整性和展示价值。
